\section{Related work}
\label{sec_related_work}

Early recommendation systems mostly rely on Collaborative Filtering (CF) \cite{schafer2007collaborative, sarwar2001item}, which are based on the idea that users with similar history will be more likely to purchase similar items. However, CF-based recommendations always have sparsity issues and cold-start problems. Therefore, some works utilize side information, like user and item attributes \cite{gong2009employing}, item contents \cite{melville2002content} to solve this issue. Among them, methods based on knowledge graph \cite{wang2020reinforced, xian2019reinforcement, zhu2020knowledge, wang2019explainable} show great advantages in the recommendation performance and explainability. 

Knowledge graph-based recommendations are roughly be categorized into embedding-based approaches and path-based approaches. Prior efforts on embedding-based knowledge graph recommendations \cite{zhang2016collaborative} always use embeddings of the knowledge graph to model the user-item interactions for recommendations. For example, exiting works \cite{huang2018improving} utilizes TranE \cite{bordes2013translating} to embed user-item interactions to integrate knowledge into the recommendation system. Similarly, \cite{grad2017graph} embeds user and item vectors into the same embedding space for recommendation. The above approaches model the relations of users and items using knowledge graph embedding methods, which achieves great improvement in model expressiveness. However, these methods are sensitive to the quality of related knowledge graphs.

%However, the embedding methods related to the knowledge graph influence the performance of the recommendation significantly. Taking the TransE as an example, it cannot handle complicated relations, like one-to-multiple relations, multiple-to-multiple relations, and so on, so its limitation is the bottleneck of TransE-based knowledge graph for the recommendation.
To this end, meta-paths \cite{sun2011pathsim}, and the connectivity of different types of nodes are introduced to recommendations, which have the ability to learn the explicit expression along each meta-path schema. In \cite{zhao2017meta}, the authors introduce Matrix Factorization (MF) and Factorization Machine (FM) to learn similarities generated by each meta-path for recommendation. \cite{han2018aspect} models rich objects and relations in knowledge graph and extracts different aspect-level similarity matrices thanks to meta-paths for the top-N recommendation. Although they achieved appealing performance, the limitation is still obvious that structured meta-path based similarity latent factors can only reflect mutual infectivity along meta-path schemas in a graph but cannot capture the specific information along particular path instances, which limits the explainability of recommendation.

More recently, injecting meta-paths as recommendation context (aggregation of meta-path instances) \cite{hu2018leveraging} was introduced for top-N recommendation. It provides fine-grained explanations based on specific instances. However, it ignores the important sequential dynamics of user-item interactions, which limits the performance of the recommendation performance and interpretability. To consider the sequential information, \cite{wang2019explainable} utilizes LSTM to model the sequential information, but it only considers path-based sequential information between users and items and ignores the importance of the user's clicked history sequences, which are highly informative to infer user's preferences. To tackle this issue, \cite{zhu2020knowledge} attempts to model the sequences of user's behaviours and path connectivity between users and items for recommendation. Nevertheless, it only considers user-item paths, which ignores the item-item intrinsic relation information and cannot learn some complex semantic information between items. Moreover, it has to pre-define some meta-paths and randomly sample meta-path instances, which involves too much human efforts and random factor. Based on the above research, we propose a reinforcement learning based path exploration for recommendation with differentiated user-item and sequential item-item instances to enhance the learning ability for explainability recommendation.



